This saturation property imposes certain restrictions on the functions $U_1,\ldots,U_6$ that specify pairwise nucleon interactions.
Suppose all the particles are confined within a region whose dimensions are of the order of the range of the nuclear forces; every pair then interacts.
If some configuration of the nucleons, including their spin orientations, makes all these pair forces attractive, the potential energy of the system is negative and proportional to $A^2$.
Its kinetic energy, by contrast, is positive and proportional to $A^{5/3}$, a lower power of $A$\footnote{The density $n$ of particles confined in the given volume is proportional to their number $A$, while the kinetic energy per particle is then proportional to $n^{2/3}$ (compare (70.1)). The total kinetic energy is therefore proportional to $AA^{2/3}$.}.
Under these conditions, a sufficiently large collection of nucleons would indeed contract into a small volume independent of $A$, and hence would not form nuclear matter.
Saturation of the nuclear forces must therefore require the absence of configurations whose negative interaction energy is proportional to $A^2$ (see problem 2).

The fact that the volume of nuclear matter is proportional to the particle number is expressed by a relation of the form
